\documentclass[../DoAn.tex]{subfiles}
\begin{document}


\section{Đặt vấn đề}
\label{section:1.1}

Trong bối cảnh kinh tế hiện đại, kĩ năng quản lý tài chính vẫn là một trong những kĩ năng quan trọng, cần thiết đối với mỗi cá nhân. Ghi chép lại chi tiêu là bước đầu trong việc quản lý tài chính, nhưng ghi chép như thế nào là thông minh và tiết kiệm thời gian.

Nhiều người vẫn đang ghi chép lại bằng thủ công như sổ tay, bảng excel tuy nhiên các phương pháp này thiếu trực quan, khó tổng hợp, mất thời gian và không hỗ trợ phân tích dài hạn. Điều này khiến người dùng không nắm rõ thói quen chi tiêu của bản thân, dẫn đến lãng phí hoặc vượt ngân sách. 

Từ đó, bài toán đặt ra là xây dựng một giải pháp hỗ trợ ghi chép, theo dõi và quản lý chi tiêu nhanh chóng, tiện lợi và chính xác. Giải pháp phù hợp không chỉ giúp người dùng kiểm soát tài chính tốt hơn mà còn hỗ trợ lập kế hoạch chi tiêu phù hợp với bản thân.

\section{Mục tiêu và phạm vi đề tài}
\label{section:1.2}
Hiện nay, nhiều ứng dụng quản lý chi tiêu cá nhân như Money Lover, Misa MoneyKeeper hay Spendee đã cung cấp các chức năng ghi chép và thống kê cơ bản. Tuy nhiên, phần lớn các ứng dụng này chủ yếu tập trung vào việc lưu trữ và hiển thị dữ liệu, trong khi khả năng phân tích hành vi chi tiêu và đưa ra gợi ý cá nhân hóa cho người dùng còn hạn chế. Ngoài ra, một số tính năng nâng cao thường yêu cầu trả phí hoặc chưa tối ưu theo thói quen sử dụng của người dùng tại Việt Nam.

Từ đó, đề tài lựa chọn phát triển một ứng dụng \textit{quản lý chi tiêu cá nhân}. Ứng dụng cho phép người dùng ghi chép thu-chi một cách linh hoạt, tiện lợi cho người dùng, theo dõi thống kê tổng quan, đồng thời đưa ra nhắc nhở cũng như gợi ý, kế hoạch chi tiêu hợp lý cho người dùng.

Phạm vi đề tài tập trung vào 3 nhóm chức năng chính: (i) Hỗ trợ ghi chép thu–chi nhanh chóng, thuận tiện và chính xác, (ii) Thống kê và báo cáo tổng quan theo thời gian, danh mục và ngân sách, (iii) Gợi ý, đề xuất kế hoạch chi tiêu phù hợp cho người dùng trong tương lai.

Đề tài không tập trung vào các chức năng nâng cao như tích hợp trực tiếp với ngân hàng, đầu tư tài chính, hay quản lý tài sản phức tạp.

\section{Bố cục đồ án}
\label{section:1.3}

Phần còn lại của báo cáo đồ án tốt nghiệp này được tổ chức như sau.

\textbf{Chương 2} trình bày quá trình phân tích yêu cầu của hệ thống, bao gồm khảo sát nhu cầu người dùng, xác định các yêu cầu chức năng, phi chức năng và các ràng buộc kỹ thuật.

\textbf{Chương 3} trình bày các công nghệ và nền tảng được sử dụng để xây dựng hệ thống.

\textbf{Chương 4} trình bày quá trình cài đặt và thực nghiệm hệ thống. Nội dung bao gồm các công nghệ được sử dụng, mô tả các chức năng đã triển khai, kết quả kiểm thử và đánh giá mức độ đáp ứng so với yêu cầu đặt ra.

\textbf{Chương 5} đưa ra kết luận và các hướng phát triển tiếp theo, tổng hợp những kết quả đạt được, chỉ ra các hạn chế và đề xuất các hướng mở rộng trong tương lai.



\end{document}
